\setchapterpreamble[u]{\margintoc}
\chapter{Le Ottimizzazioni}
\labch{Ottimizzazioni}

\section{Size}
Negli Euler Tour Trees, la struttura deve gestire sequenze 
dinamiche dove la posizione dei nodi cambia continuamente a
causa di cut e link. Per trasformare i RB-Trees una struttra
dati dinamica abbiamo bisongo di eliminare il concetto
di chiavi statiche introducendone uno che gestiche le chiavi 
in maniera implicita.\\
L'idea fondamentale è che la posizione di un nodo nella 
sequenza non è determinata da un valore fisso, ma dalla 
sua posizione relativa all'interno dell'ETT.
Per gestire efficientemente questa informazione, 
ogni nodo v memorizza un attributo aggiuntivo denominato 
\textbf{size} (tecnica formalizzata da Cormen in \cite{CLRS09})

\begin{definition}
    $size(v) = 1 + size(v.left) +size(v.right)$ 
\end{definition}

\subsubsection{Perche la size è fondamentale?}
L'introduzione della size nei Red-Black Tree permette di:
\begin{itemize}
    \item {Calcolare la posizione di un nodo nella sequenza risalendo il RB-Trees verso la radice.
        
        \begin{algorithm}[H]
            \caption{Posizione del nodo nella sequenza}
            \label{alg:find_i}
            \SetAlgoLined
            \DontPrintSemicolon
            
            \KwIn{Nodo $n$}
            \KwOut{Posizione $i$ del nodo nell'ETT}
            
            \BlankLine 
            $i \leftarrow size(n.left) + 1$\;

            \While{$n.parent \neq \text{NULL}$}{
                $p \leftarrow n.parent$\;
                
                \If{$n == p.right$}{
                    $i \leftarrow i + size(p.left) + 1$\;
                }
        
                $n \leftarrow p$\;
            }

            \Return $i$
        \end{algorithm}
    }   

    \item {Individuare l'elemento in i-esima posizione discendendo nel RB-Trees basandomi sulla size.
        
    \begin{algorithm}[H]
        \caption{Nodo in posizione i-esima}
        \label{alg:find_i-esimo}
        \SetAlgoLined
        \DontPrintSemicolon
        
        \KwIn{Radice dell'albero $r$, posizione desiderata $i$}
        \KwOut{Nodo $n$ in posizione $i$ nella sequenza}
        
        \BlankLine 
        $n \leftarrow r$\;
        
        \While{$n \neq \text{NULL}$}{
            
            \If{$i == size(n.left) + 1$}{
                \Return $n$ 
            }
            \ElseIf{$i < size(n.left) + 1$}{
                $n \leftarrow n.left$
            }
            \Else{
                $i \leftarrow i - size(n.left) + 1$\;
                $n \leftarrow n.right$ 
            }
        }
        
        \Return NULL \tcp*{Posizione non valida}
    \end{algorithm}
    }
\end{itemize}


\subsubsection{La size influenza i costi?}
L'introduzione dell'attributo \textit{size} in ogni nodo non altera la complessità temporale delle operazioni, che rimangono nell'ordine di $O(\log n)$.

\paragraph{Rotazioni}
Le rotazioni (semplici o doppie) sono operazioni locali che coinvolgono un numero costante di nodi. Poiché la \textit{size} di un nodo è
\[ size(n) = size(n.left) + size(n.right) + 1 \]
il ricalcolo del valore richiede tempo costante $O(1)$. Nelle rotazioni semplici \ref{fig:rotazioneSS}, gli unici nodi soggetti ad aggiornamento sono \textit{u} e \textit{v}; nelle rotazioni doppie \ref{fig:rotazioneSD}, l'aggiornamento interessa i nodi \textit{v, z} e \textit{w}. Di conseguenza, il costo aggiuntivo per l'aggiornamento delle size durante il ribilanciamento è $O(1)$.

\paragraph{Operazioni di Insert e Delete} 
Durante le fasi di inserimento o rimozione, il cammino percorso dalla radice verso il nodo target ha lunghezza proporzionale all'altezza dell'albero che sappiamo essere  $O(\log n)$. 
\begin{itemize}
    \item \textbf{Insert}: Ogni nodo lungo il cammino dal nodo inserito fino alla radice vede la propria \textit{size} incrementata di un'unità.
    \item \textbf{Delete}: Ogni nodo lungo il cammino del padre del nodo rimosso fino alla radice vede la propria \textit{size} decrementata di un'unità.
\end{itemize}
Tali aggiornamenti dunque non alterano il costo delle operazioni che rimane $O(\log n)$.

\paragraph{Operazioni di Join e Split} 
Nell'operazione di \texttt{Join}, il valore \textit{size} 
del nuovo nodo di giunzione $i$ viene calcolato come 
$size(i) = size(x)+ size(S_{min}) + 1 $, dove $S_{min}$ è 
l'albero con rank minore tra $S_1$ e $S_2$. Successivamente, 
risalendo il cammino da $i$ verso la radice dell'albero , 
si incrementano le \textit{size} dei nodi sul cammino di una 
quantità pari a $(size(S_{min}) + 1)$. Poiché la risalita è 
limitata dall'altezza dell'albero, il costo è $O(\log n)$.

Siccome l'operazione di \texttt{Split} è basata sull'operazioni 
di \texttt{Join}, e  il costo di quest' ultime rimane invariato 
anche la complessità dello \texttt{Split} resta invariata.

\subsection{Operazioni}

\subsubsection{isConnected}
Per eseguire l'operazione \textbf{isConnected(u,v)} bisongna verificare
che \textit{u} e \textit{v} siano nello stesso RB-Trees 
per farlo basta risalire l'albero prima da \textit{u} e poi da \textit{v} per verificare 
che la radice sia la medesima

\begin{algorithm}[H]
    \caption{isConnected}
    \label{alg:isConnectedEfficent}
    \SetAlgoLined
    \DontPrintSemicolon
    
    \KwIn{Nodo $u$, nodo $v$}
    \KwOut{True o False}
    
    \BlankLine 
    $r \leftarrow \text{NULL}$\;
    $p \leftarrow u$ \tcp*{Risale da u}\;
    
    \While{$p.parent \neq \text{NULL}$}{
        $p \leftarrow p.parent$\;
    }
    $r \leftarrow p$

    $p \leftarrow v$ \tcp*{Risale da v}\;
    \While{$p.parent \neq \text{NULL}$}{
            $p \leftarrow p.parent$\;
    }
            
    \Return $r == p$
    
\end{algorithm}

\subsubsection{Ottimizzazione tramite Tabelle Hash}

Per garantire che le operazioni di \textbf{link}, \textbf{cut} e \textbf{rerooting} mantengano una complessità temporale logaritmica $O(\log n)$, è fondamentale eliminare la dipendenza dalla scansione lineare del ciclo euleriano. Per risolvere tale problema introduciamo la struttura dati \textit{Hash Map} agisce come un dizionario che mappa chiavi univoche a riferimenti di memoria, permettendo operazioni di ricerca in tempo costante $O(1)$.
Il principio di funzionamento si basa su una funzione di hash $h: U \to \{0, \dots, m-1\}$ che mappa un universo di chiavi $U$ in un array di dimensione $m$. Per prevenire il degradamento delle prestazioni dovuto alle collisioni (quando $h(u) = h(v)$ per $u \neq v$), si adotta una famiglia di hash universale $\mathcal{H}$, che limita la probabilità di collisione tra due chiavi distinte a $1/m$.
Nella struttura degli Euler Tour Trees, l'integrazione delle tabelle hash permette di passare rapidamente dall'albero originale alla posizione fisica nel RB-Trees tramite due mappature distinte:
\begin{itemize}
    \item \textbf{mapNode: } Associa ogni nodo originale $v \in V$ all'indirizzo di memoria di un nodo nell'albero di ricerca che rappresenta un arco uscente da $v$ (ad esempio, l'arco $(v, u)$). Questa mappa è essenziale per l'operazione di \textbf{rerooting}, consentendo di individuare rapidamente il primo arco dal quale dovra partire il tour.
    \[{mapNode}[F] \to \&(F, E)\].3
    \item \textbf{mapEdge: } Associa una coppia ordinata di nodi $(u, v)$, che rappresenta un arco diretto univoco del tour, al corrispondente nodo nel RB-Trees. Questa mappatura trasforma la ricerca di un arco da una scansione lineare $O(n)$ a un accesso diretto $O(1)$, rendendo efficienti le operazioni di \textit{split} e \textit{join} necessarie per il \textbf{cut} e il \textbf{link}.
\end{itemize}


\subsubsection{Rerooting}

Per ottimizzare l'operazione ausiliaria di \textbf{rerooting}$(T, v)$, rendendola computazionalmente efficiente, si sfruttano le proprietà degli alberi binari di ricerca bilanciati (Red-Black Trees) in combinazione con una tabella hash per l'accesso diretto ai nodi. Tale approccio permette di ridurre la complessità dell'operazione da $O(n)$ a $O(\log n)$.

La procedura si articola nelle seguenti fasi:

\begin{enumerate}
    \item \textbf{Accesso alla sequenza:} Si interroga la hash map $\textit{mapNode}[v]$ per ottenere il puntatore al nodo dell'RB-Tree che memorizza un arco uscente da $v$, denotato come $e_{start} = (v, c)$. Questa operazione garantisce la localizzazione del punto di taglio in tempo costante $O(1)$.
    
    \item \textbf{Decomposizione (Split):} Si esegue l'operazione di $\text{split}(T, e_{start})$ sulla sequenza originale. Tale operazione scinde l'albero $T$ in tre componenti distinte:
    \begin{itemize}
        \item $T_L$: sottosequenza contenente tutti gli archi che precedono $e_{start}$ nel tour corrente;
        \item $e_{start}$: il nodo stesso identificato come punto di rotazione;
        \item $T_R$: sottosequenza contenente gli archi successivi a $e_{start}$.
    \end{itemize}

    \item \textbf{Rimozione della vecchia apertura:} Si estrae l'elemento situato all'estrema destra di $T_L$. Tale elemento, che denotiamo con $P$, rappresenta l'arco di inizio del tour originale.

    \item \textbf{Ricomposizione della sequenza (Join):} Si procede alla creazione della nuova sequenza $T_{new}$ invertendo l'ordine relativo dei blocchi estratti. Utilizzando $P$ come elemento di giunzione, si esegue:
    \[ T_{new} = \text{join}(T_R, P, T_L) \]

    \item \textbf{Finalizzazione del ciclo:} Per garantire la coerenza della struttura euleriana e ristabilire il nodo $v$ come radice si inserisce $e_{start}$ come elemento minimo (all'inizio) della sequenza $T_{new}$.

\end{enumerate}

\begin{marginfigure}
    \includegraphics{rerootingEfficente.png}
    \caption[Titolo breve]{Procedura di rerooting}
    \labfig{rerootingEfficente}
\end{marginfigure}

\subsubsection{Link}

L'operazione di \textbf{link}$(u, v)$ permette di unire due alberi distinti attraverso l'introduzione di un nuovo arco $(u, v)$. Per garantire un'efficienza logaritmica, l'operazione viene implementata sfruttando le struttre dati precedentemente introdotte e l'operazione ausiliaria di \textit{rerooting}.

La procedura si articola nelle seguenti fasi:

\begin{enumerate}
    \item \textbf{Normalizzazione delle radici:} Si eseguono le operazioni $\text{rerooting}(T_1, u)$ e $\text{rerooting}(T_2, v)$ sui rispettivi Euler Tour Trees. Questa fase assicura che il tour di $T_1$ inizi nel nodo $u$ e quello di $T_2$ inizi nel nodo $v$, predisponendo le sequenze per la giunzione lineare.
    
    \item \textbf{Composizione della sequenza (Join):} Si crea un nuovo nodo dell'RB-Tree, denotato come $P = (u, v)$, che funge da arco di giunzione tra i due alberi. Si procede alla concatenazione delle sottosequenze tramite l'operazione di \textit{join}:
    \[ T_{new} = \text{join}(T_1, P, T_2) \]
    Il riferimento al nuovo arco $P$ viene contestualmente inserito nelle hash map $\textit{mapEdge}$ e solo se necessario in $\textit{mapNode}$ per consentire accessi futuri in tempo costante.

    \item \textbf{Chiusura del ciclo euleriano:} Per ristabilire la proprietà di ciclo della struttura risultante e mantenere $u$ come radice complessiva, è necessario introdurre un arco di ritorno. Si crea un nodo $P' = (v, u)$ e lo si inserisce come elemento massimo (ovvero in ultima posizione) della sequenza $T_{new}$, anche l'arco $P'$ viene registrato nella struttura di supporto $\textit{mapEdge}$.
\end{enumerate}

\begin{marginfigure}
    \includegraphics{linkEfficente.png}
    \caption[Titolo breve]{Procedura di link}
    \labfig{linkEfficente}
\end{marginfigure}

\subsubsection{Cut}
L'operazione di \textbf{cut}$(u, v)$ permette di separare un albero in due alberi distinti attraverso la cancellazione di un arco $(u, v)$. Per garantire un'efficienza logaritmica, l'operazione viene implementata sfruttando le struttre dati precedentemente introdotte e l'operazione ausiliaria di \textit{rerooting}.

La procedura si articola nelle seguenti fasi:

\begin{enumerate}
    \item \textbf{Primo accesso alla sequenza:}Si interroga la hash map $\textit{mapedge}[u,v]$ per ottenere il puntatore al nodo dell'RB-Tree che memorizza un arco $e_{del}=(u,v)$. Questa operazione garantisce la localizzazione del punto di taglio in tempo costante $O(1)$.
    \item \textbf{Prima decomposizione (Split):} Si esegue l'operazione di $\text{split}(T, e_{del})$ sulla sequenza originale. Tale operazione scinde l'albero $T$ in tre componenti distinte:
    \begin{itemize}
        \item $T_L$: sottosequenza contenente tutti gli archi che precedono $e_{del}$ nel tour corrente;
        \item $e_{del}$: il nodo stesso identificato come punto di rottura;
        \item $T_R$: sottosequenza contenente gli archi successivi a $e_{del}$.
    \end{itemize}
    \item \textbf{Secondo accesso alla sequenza:}Si interroga la hash map $\textit{mapedge}[v,u]$ per ottenere il puntatore al nodo dell'RB-Tree che memorizza un arco $e'_{del}=(v,u)$.
    \item \textbf{Seconda decomposizione (Split):}Si esegue l'operazione di $\text{split}(T_R, e'_{del})$. Tale operazione scinde l'albero $T_L$ in tre componenti distinte:
    \begin{itemize}
        \item $T'_L$: sottosequenza contenente tutti gli archi successivi a $e_{del}$ e precedenti a $'e_{del}$ nel tour originale;
        \item $e'_{del}$: il nodo stesso identificato come punto di rottura;
        \item $T'_R$: sottosequenza contenente gli archi successivi a $e'_{del}$.
    \end{itemize}
    \item \textbf{Rimozione del gancio:} Si estrae l'elemento situato all'estrema destra di $T'_R$, che denotiamo con $P$.

    \item \textbf{Ricomposizione della sequenza (Join):} Si procede alla ricomposizione della sequenza utilizzando $P$ come elemento di giunzione, si esegue:
    \[ T_{new} = \text{join}(T_L, P, T'_R) \]
\end{enumerate}
In questo modo abbiamo otternuo $T'_L$ che rappresenta il tour euleriana del sottoalbero radicato in $v$ e $T_{new}$ che rappresenta l'albero originale senza il sottoalbero $T'_L$

\begin{marginfigure}
    \includegraphics{cutEfficente.png}
    \caption[Titolo breve]{Procedura di cut}
    \labfig{cutEfficente}
\end{marginfigure}